\clearpage
\thispagestyle{plain}

\begin{center}
    \textbf{\Large Resumen}
\end{center}

\small
Este trabajo desarrolla un pipeline computacional para caracterizar instancias de
\acrfull{sdp} mediante \acrfull{isa}, comparar el rendimiento de solvers
representativos y preparar la generación dirigida de nuevas instancias sintéticas.
La metodología integra la recopilación de instancias desde bibliotecas de benchmarks,
su normalización al formato SDPA sparse, la ejecución automatizada de solvers, la
extracción de features estructurales, la generación de metadata compatible con ISA y
la construcción preliminar del espacio de instancias. Además, se implementa una
aplicación que permite inspeccionar visualmente los artefactos generados, proyectar
metadata de prueba mediante \texttt{exploreIS} e identificar regiones vacías del
espacio como objetivos para una etapa posterior basada en algoritmos genéticos.

Al cierre de esta etapa, el sistema permite ejecutar el flujo experimental de extremo
a extremo y producir artefactos interpretables, incluyendo coordenadas del espacio de
instancias, footprints preliminares de desempeño y puntos candidatos en zonas
subrepresentadas. Los resultados obtenidos no constituyen aún una caracterización
definitiva del rendimiento de los solvers, pero entregan evidencia inicial de que las
instancias SDP presentan variación estructural observable y que dicha variación puede
ser analizada mediante ISA para orientar tanto la comparación algorítmica como la
generación evolutiva de nuevas instancias.

\vspace{0.45cm}

\begin{center}
    \textbf{\Large Abstract}
\end{center}

This work develops a computational pipeline to characterize \acrfull{sdp} instances
through \acrfull{isa}, compare the performance of representative solvers, and prepare
the guided generation of new synthetic instances. The methodology integrates the
collection of benchmark instances, their normalization to the SDPA sparse format,
automated solver execution, structural feature extraction, generation of
ISA-compatible metadata, and the preliminary construction of the instance space. In
addition, an application is implemented to visually inspect the generated artifacts,
project test metadata through \texttt{exploreIS}, and identify empty regions of the
space as targets for a later stage based on genetic algorithms.

At the end of this stage, the system can execute the experimental workflow from end to
end and produce interpretable artifacts, including instance-space coordinates,
preliminary performance footprints, and candidate points in underrepresented regions.
The obtained results do not yet provide a definitive characterization of solver
performance, but they offer initial evidence that SDP instances exhibit observable
structural variation and that this variation can be analyzed through ISA to guide both
algorithm comparison and the evolutionary generation of new instances.

\normalsize
